\documentclass[11pt,a4paper]{article}
\usepackage{fontspec}
\usepackage{xeCJK}
\setCJKmainfont{STSong}
\setCJKsansfont{STHeiti}
\usepackage{amsmath,amssymb,amsthm}
\usepackage{mathtools}
\usepackage{bm}
\usepackage{booktabs}
\usepackage{array}
\usepackage{enumitem}
\usepackage{geometry}
\usepackage{xcolor}
\usepackage{tcolorbox}
\usepackage{algorithm}
\usepackage{algpseudocode}
\usepackage{pifont}
\usepackage{hyperref}

\geometry{margin=2.5cm}

\newcommand{\E}{\mathbb{E}}
\newcommand{\R}{\mathbb{R}}
\newcommand{\KL}{\mathrm{KL}}
\newcommand{\N}{\mathcal{N}}
\newcommand{\PC}{\mathrm{PCGrad}}
\newcommand{\vb}{v_{\mathrm{base}}}
\newcommand{\vfused}{v_{\mathrm{fused}}}
\newcommand{\inner}[2]{\langle #1,\, #2\rangle}
\newcommand{\norm}[1]{\lVert #1\rVert}
\newcommand{\proj}[2]{\mathrm{proj}_{#2}\!\left(#1\right)}
\DeclareMathOperator{\sign}{sign}
\DeclareMathOperator{\Var}{Var}
\DeclareMathOperator*{\softmax}{softmax}

\theoremstyle{plain}
\newtheorem{theorem}{Theorem}
\newtheorem{proposition}{Proposition}
\newtheorem{lemma}{Lemma}
\newtheorem{corollary}{Corollary}

\theoremstyle{definition}
\newtheorem{definition}{Definition}
\newtheorem{remark}{Remark}
\newtheorem{example}{Example}

\title{\textbf{Ensemble Blend Modes: Definitions and Theory}\\[0.3em]
\large \texttt{ensemble-eval} 多 teacher 聚合算法（weighted / PCGrad 家族 / TIES）的含义与理论分析}
\author{Flow Factory}
\date{}

\begin{document}
\maketitle

\begin{abstract}
本文系统给出 \texttt{ensemble-eval} trainer 在每个去噪步对 $K$ 个 LoRA checkpoint 的 \texttt{noise\_pred} 进行融合时所用的全部聚合算法（\texttt{ensemble\_blend\_mode}）的严格定义与理论分析，并解释 \texttt{scripts/compare\_ensemble\_methods.py} 默认对比的那一组方法为何如此选择。核心结论：
\begin{enumerate}[leftmargin=1.6em]
\item \textbf{无冲突即退化}（命题~\ref{prop:noconflict}）：任意 PCGrad 投影（global / channelwise / normalized）在没有符号冲突时都退化为加权和 \texttt{weighted}。单 checkpoint 时三者恒等于 \texttt{weighted}。
\item \textbf{全速度坍缩}（命题~\ref{prop:collapse}）：teacher 的 full velocity $v_i$ 被共享的 base 方向主导（近平行），冲突率趋零，故 \texttt{pcgrad}/\texttt{pcgrad\_channelwise}/\texttt{pcgrad\_normalized} 与 \texttt{weighted} 数值上几乎相同。
\item \textbf{base 锚定暴露冲突}（命题~\ref{prop:residual}）：在残差 $\tau_i=v_i-\vb$ 上操作移除了共享共模分量，task delta 之间冲突频繁，故 \texttt{pcgrad\_residual} 系列与 \texttt{ties} 才与基线实质不同。
\item \textbf{TIES 的退化与去毒}（命题~\ref{prop:ties}）：无符号分歧时 TIES 退化为加权平均；有分歧时通过 sign election 丢弃逆向 teacher，等价于一种逐元素的鲁棒聚合。
\end{enumerate}
据此，默认对比集合 $=\{$\texttt{weighted}（基线）$\}\cup\{$\texttt{pcgrad\_residual} 系列 $+$ \texttt{ties}$\}\cup\{$KL/$1/$KL 动态加权$\}$，外加 base 与各 teacher 的单独 baseline。
\end{abstract}

\tableofcontents
\newpage

%% ============================================================
\section{设定与符号}
%% ============================================================

\texttt{ensemble-eval} trainer（\texttt{src/flow\_factory/trainers/ensemble\_eval/}）在每个去噪步把 $K$ 个 checkpoint 的 \texttt{noise\_pred} 融合为单个 $\vfused$，再做一次 \texttt{scheduler.step}。\emph{输出空间} 融合算子的入口是 \texttt{ensemble\_forward\_step}（\texttt{common.py}），其分发逻辑为：
\[
\text{vector}\in\{\underbrace{v_i}_{\text{full velocity}},\ \underbrace{\tau_i=v_i-\vb}_{\text{residual delta}}\}
\ \times\
\text{projection}\in\{\text{global},\ \text{channelwise},\ \text{normalized}\}
\ \cup\ \{\texttt{weighted},\ \texttt{ties},\ \texttt{ties\_channelwise}\}.
\]
此外 \texttt{weight\_merge}（\S\ref{sec:weight_merge}）在 \emph{参数空间} 一次性融合（不经 per-step 路径），作为 model-soup baseline。

\paragraph{符号约定。} 所有量均为 \emph{per-sample}（固定 batch 样本 $b$ 与时间步 $k$，省略上下标）：
\begin{center}
\begin{tabular}{cl}
\toprule
\textbf{符号} & \textbf{含义} \\
\midrule
$K$ & ensemble 中 checkpoint（teacher）数 \\
$d=C\times H\times W$ & latent 维度 \\
$\vb$ & base（pretrained，LoRA 关闭）的 velocity（\texttt{use\_ref\_parameters}） \\
$v_i$ & teacher $i$ 的 velocity（\texttt{noise\_pred}，\texttt{use\_named\_parameters}） \\
$\tau_i=v_i-\vb$ & teacher $i$ 的 task vector（residual delta） \\
$\pi_i\ge0,\ \sum_i\pi_i=1$ & 静态 checkpoint 权重（\texttt{checkpoint\_weights}，默认均匀 $1/K$） \\
$w_i\ge0$ & 实际融合权重（静态或 per-sample 动态，\S\ref{sec:weighting}） \\
$\varepsilon$ & \texttt{pcgrad\_eps}，数值下限（默认 $10^{-8}$） \\
\bottomrule
\end{tabular}
\end{center}

\noindent 投影记号：对向量 $a$ 在 $b$ 方向的投影 $\proj{a}{b}=\dfrac{\inner{a}{b}}{\norm{b}^2}\,b$。

%% ============================================================
\section{Weighted：线性基线}
%% ============================================================

\begin{definition}[\texttt{weighted}]
\label{def:weighted}
线性加权和（\texttt{common.py}, full-velocity 分支）：
\begin{equation}
\vfused^{\,\mathrm{wsum}} \;=\; \sum_{i=1}^{K} w_i\, v_i.
\label{eq:weighted}
\end{equation}
\end{definition}

这是最简单的聚合：各 teacher 速度的凸组合（$w_i$ 归一化时）。它不做任何冲突检测，是其余所有方法的参照系；下文反复出现的结论是"某模式在某条件下退化为~\eqref{eq:weighted}"。从 RL/蒸馏视角，\eqref{eq:weighted} 等价于把各 teacher 的转移均值做加权平均（见 \texttt{teacher\_aggregation\_loss\_formulations.tex}），即 \texttt{average}/\texttt{sum} 聚合的推理期对应物。

%% ============================================================
\section{PCGrad 家族：冲突投影}
%% ============================================================

PCGrad（Projecting Conflicting Gradients，Yu et al.\ 2020）原用于多任务梯度去冲突。这里把"梯度"换成 teacher 的速度/残差向量。

\subsection{Global PCGrad}

\begin{definition}[\texttt{pcgrad}，global]
\label{def:pcgrad}
输入 $\{g_i\}$（full 模式 $g_i=w_iv_i$，残差模式 $g_i=w_i\tau_i$，见 \texttt{\_blend\_velocity\_set}）。初始化 $\mathrm{pc}_i\gets g_i$；对每个 $i$，按固定（含可复现随机）顺序遍历 $j\neq i$：
\begin{equation}
\text{若 } \inner{\mathrm{pc}_i}{g_j}<0:\quad
\mathrm{pc}_i \;\gets\; \mathrm{pc}_i - \frac{\inner{\mathrm{pc}_i}{g_j}}{\norm{g_j}^2}\,g_j
\;=\; \mathrm{pc}_i - \proj{\mathrm{pc}_i}{g_j},
\label{eq:pcgrad_step}
\end{equation}
投影方向始终用 \emph{原始} $g_j$。输出 $\vfused=\sum_i\mathrm{pc}_i$。冲突判定 $\inner{\cdot}{\cdot}<0$ 在全部 $d$ 维上做一次内积。
\end{definition}

几何含义：当 teacher $i$ 与 $j$ 方向相悖（夹角 $>90^\circ$）时，把 $i$ 中与 $j$ 相抵的分量减去，只保留非冲突部分，再求和。其作用是\emph{抑制相互抵消}，让融合方向更"协调"。

\subsection{Channelwise PCGrad}

\begin{definition}[\texttt{pcgrad\_channelwise}]
\label{def:channelwise}
与~\eqref{eq:pcgrad_step} 相同，但内积与冲突判定在更细粒度上独立进行（\texttt{pcgrad\_blend\_noise\_preds\_channelwise}）：
\begin{itemize}[leftmargin=1.6em]
\item 4D $(B,C,H,W)$：\emph{逐 channel}，内积在 $H\times W$ 上，每个 channel 独立决定是否投影；
\item 3D $(B,\text{seq},\text{feat})$：\emph{逐 token}，内积在 feat 上。
\end{itemize}
\end{definition}

直觉：global 用一个标量内积概括整张图的冲突，可能"一票否决"；channelwise 允许部分 channel/token 冲突、部分不冲突，去冲突更局部、更激进（冲突触发面更大）。

\subsection{Normalized PCGrad}

\begin{definition}[\texttt{pcgrad\_normalized}]
\label{def:normalized}
把每个向量拆成幅度 $n_i=\norm{v_i}$ 与单位方向 $u_i=v_i/n_i$（\texttt{pcgrad\_blend\_noise\_preds\_normalized}）。PCGrad 在 \emph{单位方向} $\{u_i\}$ 上运行~\eqref{eq:pcgrad_step} 得投影方向 $p_i$，再恢复幅度与权重重组：
\begin{equation}
\vfused \;=\; \sum_{i=1}^{K} w_i\, n_i\, p_i.
\label{eq:normalized}
\end{equation}
\end{definition}

动机：global 的投影系数 $\inner{\mathrm{pc}_i}{g_j}/\norm{g_j}^2$ 受 $g_j$ 幅度影响，高范数 teacher 会主导投影。归一化后冲突几何 $\inner{u_i}{u_j}$（即 cosine）幅度无关，避免"大嗓门"teacher 支配方向；幅度只在最后线性重组时加回。

\subsection{Residual（base-anchored）变体}

\begin{definition}[\texttt{pcgrad\_residual} 系列]
\label{def:residual}
先做一次 base 前向得 $\vb$（\texttt{use\_ref\_parameters}），各 teacher 残差 $\tau_i=v_i-\vb$；对 $\{\tau_i\}$（而非 $\{v_i\}$）施加上述三种投影之一得 $\{\tau_i^{\PC}\}$，再加回基线（\texttt{\_pcgrad\_residual\_blend}）：
\begin{equation}
\vfused \;=\; \vb \;+\; \underbrace{\textstyle\sum_i w_i\,\tau_i^{\PC}}_{\text{combined delta}}.
\label{eq:residual}
\end{equation}
对应 \texttt{pcgrad\_residual}（global）/ \texttt{\_channelwise} / \texttt{\_normalized}。代价：每步多一次 base 前向。
\end{definition}

%% ============================================================
\section{TIES：符号选举 $+$ 不相交合并}
%% ============================================================

\begin{definition}[\texttt{ties}]
\label{def:ties}
base-anchored。在 task vectors $\{\tau_i\}$ 上做 TIES-merging（Yadav et al.\ 2023）的两步（\texttt{ties\_blend\_deltas}）：
\begin{enumerate}[leftmargin=1.6em]
\item（可选）\textbf{trim}：每个 $\tau_i$ 按幅度保留 top-$\rho$ 比例的元素（\texttt{ties\_density}$=\rho$，默认 $1.0$ 不裁剪），其余置零；
\item \textbf{sign election}：逐元素选举主导符号
\begin{equation}
\gamma \;=\; \sign\Bigl(\textstyle\sum_i w_i\,\tau_i\Bigr)\in\{-1,0,+1\}^d;
\label{eq:ties_sign}
\end{equation}
\item \textbf{disjoint merge}：仅对与 $\gamma$ 同号的 teacher 做加权平均
\begin{equation}
\tau_{\mathrm{merged}}
=\frac{\sum_i w_i\,\tau_i\,\mathbf{1}[\sign(\tau_i)=\gamma]}{\sum_i w_i\,\mathbf{1}[\sign(\tau_i)=\gamma]}
\quad(\text{分母为 }0\text{ 处取 }0),
\label{eq:ties_merge}
\end{equation}
\end{enumerate}
返回 $\vfused=\vb+\tau_{\mathrm{merged}}$。
\end{definition}

直觉：逐元素地"投票"决定该坐标应往哪个方向修正，再只让赞成方贡献，逆向 teacher 被直接剔除（而非像 PCGrad 那样按投影削减）。这是一种比 PCGrad 更"硬"的冲突消解。

\begin{definition}[\texttt{ties\_channelwise}]
\label{def:ties_channelwise}
TIES 的 channel 粒度变体（group $=$ \texttt{dim 1}，与 channelwise PCGrad 同一分组；要求 ndim $\ge 3$）。先按 feature 维求每个 teacher 的 \emph{每通道净质量} $m_{i,c}=\sum_{\text{feat}}\tau_i[c,\cdot]$，据此 \emph{每通道} 选举一个符号
\begin{equation}
\gamma_c \;=\; \sign\Bigl(\textstyle\sum_i w_i\,m_{i,c}\Bigr);
\label{eq:ties_cw_sign}
\end{equation}
再做 \emph{整通道} 的 disjoint 加权平均——仅当 teacher 的该通道净符号 $\sign(m_{i,c})=\gamma_c$ 时，其 \emph{整条通道} 才参与
\begin{equation}
\tau_{\mathrm{merged}}[c,\cdot]
=\frac{\sum_i w_i\,\tau_i[c,\cdot]\,\mathbf{1}[\sign(m_{i,c})=\gamma_c]}{\sum_i w_i\,\mathbf{1}[\sign(m_{i,c})=\gamma_c]}.
\label{eq:ties_cw_merge}
\end{equation}
返回 $\vfused=\vb+\tau_{\mathrm{merged}}$。与逐元素 \texttt{ties} 相比，决策与动作都在通道粒度（更粗、更"成块"），对应 PCGrad 的 global$\to$channelwise 类比。
\end{definition}

%% ============================================================
\section{Weight Merge：权重空间融合（model soup）}
\label{sec:weight_merge}
%% ============================================================

上述所有模式都在 \emph{输出空间}（per-step velocity）上融合，每步需 $K$ 次（residual/ties 为 $K{+}1$ 次）前向。\texttt{weight\_merge} 则在 \emph{参数空间} 上融合，且只融合一次。

\begin{definition}[\texttt{weight\_merge}]
\label{def:weight_merge}
设 teacher $i$ 的 LoRA 参数为 $\theta_i$（与 base 架构对齐的张量列表）。一次性做权重空间加权平均，得到单个合并 LoRA：
\begin{equation}
\theta_{\mathrm{merged}} \;=\; \sum_{i=1}^{K} w_i\,\theta_i,\qquad w_i=\pi_i\ (\text{静态 \texttt{checkpoint\_weights}，默认 }1/K),
\label{eq:weight_merge}
\end{equation}
随后用 $\theta_{\mathrm{merged}}$ 做 \emph{普通单模型} 推理（每步仅 1 次前向，无 per-step blend、无 $\vb$）。实现见 \texttt{build\_merged\_lora\_snapshot}：对齐的 teacher 快照逐槽 fp32 加权求和，写入合并快照后经 \texttt{use\_named\_parameters} 单模型推理。
\end{definition}

\begin{remark}[LoRA 双线性：近似而非精确 delta merge]
\label{rem:bilinear}
LoRA 的有效增量为 $\Delta W_i=B_iA_i$（双线性）。\eqref{eq:weight_merge} 对 $A_i,B_i$ \emph{分别} 平均，故
\[
\Bigl(\textstyle\sum_i w_iB_i\Bigr)\Bigl(\sum_i w_iA_i\Bigr)\ \neq\ \sum_i w_i\,B_iA_i
\]
即 \texttt{weight\_merge} 复现的是左式（廉价的 LoRA soup），\emph{不等于} 右式的精确 task-arithmetic delta merge。这是有意为之的简单基线。又因无 per-step $\vb$，KL 动态加权不适用，\texttt{ensemble\_blend\_weighting} 必须为 \texttt{uniform}。
\end{remark}

\begin{remark}[系数性质与成本]
\label{rem:wm_props}
权重 $\theta_{\mathrm{merged}}$ 是 teacher 参数的凸组合（$w_i\ge0,\sum w_i=1$），但由于扩散映射对参数高度非线性，\emph{输出} $\vfused$ 一般不是各 teacher 输出 $v_i$ 的凸组合（与 \texttt{weighted} 的输出空间凸组合不同）。成本上最省：合并一次后每步单前向，远低于 PCGrad/TIES 的 $K$/$K{+}1$ 次。
\end{remark}

%% ============================================================
\section{核心理论：退化关系与 base 锚定的必要性}
%% ============================================================

本节给出"为何默认对比只保留 \texttt{weighted} $+$ residual 系列 $+$ \texttt{ties}"的严格依据。

\begin{lemma}[投影对正缩放的齐次性]
\label{lem:homog}
固定遍历顺序与冲突判定，~\eqref{eq:pcgrad_step} 对输入正缩放齐次：若 $g_i\mapsto c\,g_i$（$c>0$ 对所有 $i$ 相同），则 $\mathrm{pc}_i\mapsto c\,\mathrm{pc}_i$。
\end{lemma}
\begin{proof}
归纳。内积 $\inner{c\,\mathrm{pc}_i}{c\,g_j}=c^2\inner{\mathrm{pc}_i}{g_j}$ 符号不变，故冲突判定一致；投影 $\proj{c\,\mathrm{pc}_i}{c\,g_j}=c\,\proj{\mathrm{pc}_i}{g_j}$，更新两边同乘 $c$。
\end{proof}

\begin{proposition}[无冲突即退化为 weighted]
\label{prop:noconflict}
若在某样本上 PCGrad 内层循环 \emph{未触发任何投影}（所有判定 $\inner{\cdot}{\cdot}\ge0$），则
\begin{itemize}[leftmargin=1.6em]
\item global / channelwise：$\mathrm{pc}_i=g_i=w_i v_i$（或 $w_i\tau_i$），故 $\vfused=\sum_i w_i v_i$，即 \texttt{weighted}（残差模式为 $\vb+\sum_i w_i\tau_i$）；
\item normalized：$p_i=u_i$，重组 $\sum_i w_i n_i u_i=\sum_i w_i v_i$，\emph{代数上} 亦等于 \texttt{weighted}。
\end{itemize}
特别地，$K=1$ 时不存在 $j\neq i$，三种投影 \emph{恒} 退化为 $w_1 v_1$。
\end{proposition}
\begin{proof}
无投影时 $\mathrm{pc}_i$ 保持初值；global/channelwise 初值即 $g_i$。normalized 初值 $p_i=u_i$，代入~\eqref{eq:normalized} 得 $\sum_i w_i n_i u_i=\sum_i w_i v_i$。$K=1$ 见定义中的单元素短路分支。
\end{proof}

\begin{remark}
命题~\ref{prop:noconflict} 解释了脚本中"单 teacher baseline"为何直接用 \texttt{blend\_mode=weighted} 配单 checkpoint：由短路分支它就是该 teacher 自身的 $v_1$。
\end{remark}

\subsection{为何全速度 PCGrad 坍缩到 weighted}

\begin{proposition}[full-velocity 坍缩]
\label{prop:collapse}
设各 teacher 由同一 base 微调，故 $v_i=\vb+\tau_i$ 且 $\norm{\tau_i}\ll\norm{\vb}$（残差远小于共享分量）。则任意两 teacher 的 full velocity 夹角余弦
\begin{equation}
\cos\angle(v_i,v_j)
=\frac{\inner{\vb+\tau_i}{\vb+\tau_j}}{\norm{\vb+\tau_i}\,\norm{\vb+\tau_j}}
= 1 - O\!\Bigl(\tfrac{\norm{\tau}^2}{\norm{\vb}^2}\Bigr)\;\to\;1,
\label{eq:collapse}
\end{equation}
即 $v_i$ 近平行、内积恒正。于是 global PCGrad 几乎不触发投影，由命题~\ref{prop:noconflict} 数值上 $\approx$ \texttt{weighted}；normalized 同理（单位方向近重合）。
\end{proposition}
\begin{proof}
展开 $\inner{\vb+\tau_i}{\vb+\tau_j}=\norm{\vb}^2+\inner{\vb}{\tau_i+\tau_j}+\inner{\tau_i}{\tau_j}$，分母 $\norm{\vb+\tau_i}=\norm{\vb}\sqrt{1+2\inner{\vb}{\tau_i}/\norm{\vb}^2+O(\norm{\tau}^2/\norm{\vb}^2)}$。令 $\delta=\norm{\tau}/\norm{\vb}\to0$，逐阶展开即得 $\cos=1-O(\delta^2)$。正内积 $\Rightarrow$ 冲突判定 $\inner{\cdot}{\cdot}<0$ 几乎不成立。
\end{proof}

\begin{corollary}
\label{cor:exclude}
\texttt{pcgrad}, \texttt{pcgrad\_channelwise}, \texttt{pcgrad\_normalized} 在 fine-tuned ensemble 上与 \texttt{weighted} 数值近乎一致，对比中信息冗余，故默认对比集合将其剔除（仅保留 \texttt{weighted} 作为这一族的代表基线）。channelwise 因逐 channel 判定冲突面更大，可能略有偏离，但共模主导对每个 channel 同样成立，定性不变。
\end{corollary}

\subsection{为何 base 锚定（残差）才有效}

\begin{proposition}[残差移除共模、暴露冲突]
\label{prop:residual}
在残差空间，融合对象是 $\tau_i=v_i-\vb$，共享分量 $\vb$ 被精确减去。task delta 之间无共模约束，夹角可自由分布于 $[0,\pi]$；当 teacher 专精不同目标时，$\inner{\tau_i}{\tau_j}<0$（方向相悖）频繁出现。因此 residual PCGrad/\texttt{ties} 的投影/选举\emph{真正触发}，融合结果与 \texttt{weighted}（残差）$\vb+\sum_i w_i\tau_i$ 实质不同。
\end{proposition}
\begin{proof}
由构造 $\tau_i$ 不含 $\vb$ 方向的强制正分量；命题~\ref{prop:collapse} 的 $\cos\to1$ 论证不再适用。冲突率由 $\{\tau_i\}$ 的真实几何决定，可显著 $>0$（经验上由 \texttt{PCGradStats} 的 conflict 率证实：full 模式 $\approx0$，residual 模式显著为正）。
\end{proof}

\begin{tcolorbox}[colback=gray!6,colframe=gray!50,title=\textbf{为何默认对比是这一组}]
\begin{itemize}[leftmargin=1.4em]
\item \texttt{weighted}：线性基线（也代表整族 full-velocity，命题~\ref{prop:collapse}）。
\item \texttt{pcgrad\_residual} / \texttt{\_channelwise} / \texttt{\_normalized}：唯一与基线实质不同的 PCGrad（命题~\ref{prop:residual}）。
\item \texttt{ties}：base-anchored 的硬冲突消解，独立于 PCGrad 的一类。
\item \texttt{pcgrad\_residual\_kl} / \texttt{\_kl\_inv}：在 residual 上叠加 per-sample 动态加权（\S\ref{sec:weighting}）。
\item baseline：\texttt{base}（$\vb$）与各 teacher 单独（$v_i$），给出融合的下界参照。
\end{itemize}
被剔除：\texttt{pcgrad} / \texttt{pcgrad\_channelwise} / \texttt{pcgrad\_normalized}（坍缩到 \texttt{weighted}）。
\end{tcolorbox}

\subsection{TIES 的退化与鲁棒性}

\begin{proposition}[TIES 退化为加权平均]
\label{prop:ties}
若所有 teacher 在某元素上符号一致（无分歧），则~\eqref{eq:ties_merge} 的指示函数恒为 $1$，
\[
\tau_{\mathrm{merged}}=\frac{\sum_i w_i\tau_i}{\sum_i w_i}=\sum_i \tilde w_i\,\tau_i,\qquad \tilde w_i=w_i/\textstyle\sum_j w_j,
\]
即归一化加权平均（$\rho=1$ 时）。存在分歧时，逆向（与 $\gamma$ 异号）teacher 在该元素被剔除，等价于逐元素的 sign-trimmed 加权平均——对"少数派强烈反向"具有鲁棒性。
\end{proposition}
\begin{proof}
指示函数恒 $1$ 时分子分母同为 $\sum_i w_i(\cdot)$，直接化简。分歧时异号项被乘零剔除，仅同号项进入加权平均。
\end{proof}

\begin{remark}[与 PCGrad 的对比]
PCGrad 对冲突做\emph{连续}的投影削减（保留正交分量）；TIES 做\emph{离散}的符号剔除（整项丢弃）。前者温和、保几何；后者激进、保符号一致性。二者都需 base 锚定才有意义（命题~\ref{prop:residual}）。
\end{remark}

%% ============================================================
\section{动态加权：KL / inverse-KL}
\label{sec:weighting}
%% ============================================================

权重 $w_i$ 可为静态 $\pi_i$（\texttt{uniform}）或 per-sample 动态。动态族以 teacher--base 的 per-step KL（在 flow matching 下 $=$ velocity-MSE）$D_i=\norm{\tau_i}^2$ 的幂为权：
\begin{equation}
w_i(\alpha)=\frac{\pi_i\,D_i^{\alpha}}{\sum_j \pi_j\,D_j^{\alpha}},\qquad
\begin{cases}
\alpha=0 & \texttt{uniform}\\
\alpha=+1 & \texttt{kl}\ (\text{放大偏离最大的 specialist，软路由})\\
\alpha=-1 & \texttt{kl\_inv}\ (\text{inverse-variance，压制偏离大者}).
\end{cases}
\label{eq:kl_weights}
\end{equation}
因需要 $\vb$，\texttt{kl}/\texttt{kl\_inv} 仅对 base-anchored 模式（residual 系列与 \texttt{ties}）有效（\texttt{\_resolve\_blend\_weights} / \texttt{\_dynamic\_kl\_weights}，否则显式报错）。

\begin{theorem}[加权与 global/channelwise 投影可交换]
\label{thm:commute}
对 global/channelwise PCGrad，引入 per-sample 标量 $w_i>0$ 等价于先用 \emph{未加权} $\{\tau_j\}$ 投影、再加权重组：$\mathrm{pc}_i=w_i\,\tau_i^{\PC}$，从而 $\vfused=\vb+\sum_i w_i\tau_i^{\PC}$。即动态加权只改变最终线性重组，不改变冲突结构。
\end{theorem}
\begin{proof}
即引理~\ref{lem:homog} 的逐 teacher 推广（不同 $w_i$）：内积 $\inner{w_i\mathrm{pc}_i}{w_j\tau_j}=w_iw_j\inner{\mathrm{pc}_i}{\tau_j}$，$w_iw_j>0$ 符号不变；投影中 $w_j$ 于分子分母相消，更新保持 $\mathrm{pc}_i=w_i\tau_i^{(s)}$。详见 \texttt{kl\_weighted\_teacher\_fusion.tex} 定理~1。
\end{proof}

\noindent normalized 与 \texttt{ties} 不满足可交换（前者把幅度剥离、$w_i\propto\norm{\tau_i}^{2\alpha}$ 改变等效幂；后者 sign election 显式用 $w_i$）。\texttt{ensemble-eval} 的 OCR/PickScore/GenEval 属 specialist 制度，理论预测 $\alpha>0$（\texttt{kl}）$\gg$ \texttt{uniform} $\gg$ $\alpha<0$（\texttt{kl\_inv}）。完整推导见 \texttt{kl\_weighted\_teacher\_fusion.tex}。

%% ============================================================
\section{速查表}
%% ============================================================

\begin{center}
\begin{tabular}{lllc}
\toprule
\textbf{mode} & \textbf{融合公式} & \textbf{冲突处理} & \textbf{base 前向} \\
\midrule
\texttt{weighted} & $\sum_i w_i v_i$ & 无 & \ding{55} \\
\texttt{pcgrad} & $\sum_i \mathrm{pc}_i(\{w_jv_j\})$ & global 投影 & \ding{55} \\
\texttt{pcgrad\_channelwise} & 同上，逐 channel/token & 细粒度投影 & \ding{55} \\
\texttt{pcgrad\_normalized} & $\sum_i w_i n_i p_i$ & 单位方向投影 & \ding{55} \\
\texttt{pcgrad\_residual} & $\vb+\sum_i w_i\tau_i^{\PC}$ & global 投影（残差） & \ding{51} \\
\texttt{pcgrad\_residual\_channelwise} & 同上，逐 channel & 细粒度（残差） & \ding{51} \\
\texttt{pcgrad\_residual\_normalized} & $\vb+\sum_i w_i n_i^\tau p_i^\tau$ & 单位方向（残差） & \ding{51} \\
\texttt{ties} & $\vb+\tau_{\mathrm{merged}}$ & 逐元素符号选举$+$剔除 & \ding{51} \\
\texttt{ties\_channelwise} & $\vb+\tau_{\mathrm{merged}}$ & 逐通道符号选举$+$整通道剔除 & \ding{51} \\
\texttt{weight\_merge} & 单模型 $f(\theta_{\mathrm{merged}})$，$\theta_{\mathrm{merged}}{=}\sum_i w_i\theta_i$ & 无（参数空间） & \ding{55} \\
\bottomrule
\end{tabular}
\end{center}

\noindent 加权选项 \texttt{ensemble\_blend\_weighting}$\in\{$\texttt{uniform}, \texttt{kl}, \texttt{kl\_inv}$\}$ 仅对带 \ding{51}（base-anchored）的行有效；\texttt{weight\_merge} 仅支持 \texttt{uniform}（无 per-step $\vb$）。注意 \texttt{weight\_merge} 在 \emph{参数空间} 融合、每步仅 1 次前向，其余模式在 \emph{输出空间} 融合。

\begin{center}
\begin{tabular}{lll}
\toprule
\textbf{模式} & \textbf{与 weighted 关系} & \textbf{默认对比} \\
\midrule
\texttt{pcgrad}/\texttt{\_channelwise}/\texttt{\_normalized} & 近似相等（命题~\ref{prop:collapse}） & 剔除 \\
\texttt{weighted} & --- & 保留（基线） \\
\texttt{pcgrad\_residual} 系列 & 实质不同（命题~\ref{prop:residual}） & 保留 \\
\texttt{ties} / \texttt{ties\_channelwise} & 实质不同（逐元素 / 逐通道） & 保留 \\
\texttt{*\_kl} / \texttt{*\_kl\_inv} & 残差 $+$ 动态重组 & 保留 \\
\texttt{weight\_merge} & 参数空间（非输出空间） & 保留（baseline） \\
\bottomrule
\end{tabular}
\end{center}

%% ============================================================
\section{结论}
%% ============================================================

\begin{enumerate}[leftmargin=1.6em]
\item 所有 PCGrad 投影在无冲突时退化为 \texttt{weighted}（命题~\ref{prop:noconflict}）；单 checkpoint 时恒等。
\item full-velocity 被共享 base 方向主导，冲突率趋零，\texttt{pcgrad}/\texttt{\_channelwise}/\texttt{\_normalized} 坍缩到 \texttt{weighted}（命题~\ref{prop:collapse}），故默认对比剔除之。
\item base 锚定（残差）移除共模、暴露真实冲突，是 PCGrad/\texttt{ties} 产生非平凡效果的前提（命题~\ref{prop:residual}）。
\item TIES 无分歧时退化为加权平均，有分歧时硬剔除逆向 teacher（命题~\ref{prop:ties}），与 PCGrad 的软投影互补。
\item 动态 KL 加权与 global/channelwise 投影可交换（定理~\ref{thm:commute}），是低成本、可叠加的正交改进；符号 $\alpha$ 由 ensemble 制度决定。
\end{enumerate}

\paragraph{一句话总结。} 在"各 teacher 由同一 base 微调"的 \texttt{ensemble-eval} 设定下，唯有在 \emph{残差} $\tau_i=v_i-\vb$ 上做冲突消解（\texttt{pcgrad\_residual} 系列、\texttt{ties}）才与简单加权和不同；全速度版本因共模主导而退化。默认对比因此聚焦于 \texttt{weighted} 基线、residual 冲突消解族、动态加权变体，以及 base/teacher 单独 baseline。

\paragraph{相关文档。}
\texttt{kl\_weighted\_teacher\_fusion.tex}（KL 动态加权的完整理论）、
\texttt{pcgrad\_gradient\_vs\_velocity\_analysis.tex}（gradient-space vs velocity-space PCGrad）、
\texttt{teacher\_aggregation\_loss\_formulations.tex}（训练期 teacher aggregation）。

\end{document}
